\section*{Capítulo \ref{ch:b1:linmaps} --- Aplicaciones lineales}

\begin{solution}{exo:b1:linmaps:1}
(1) \omterm{def:b1:linmaps:def}{Lineal}: las \omterm{prop:b1:vspaces:coordinates}{coordenadas} son expresiones \omterm{def:b1:linmaps:def}{lineales}. (2) No
\omterm{def:b1:linmaps:def}{lineal}: $u(2(1,1)) = 4 \neq 2 = 2u(1,1)$. (3) \omterm{def:b1:linmaps:def}{Lineal}: lo son
derivar y multiplicar por $X$, y las sumas de \omterm{def:b1:linmaps:def}{aplicaciones lineales}
también. (4) \omterm{def:b1:linmaps:def}{Lineal}: la evaluación respeta las operaciones punto a
punto.
\end{solution}

\begin{solution}{exo:b1:linmaps:2}
\omterm{def:b1:linmaps:kerim}{Núcleo}: resuélvase $x + y - z = 0$, $2x + y + z = 0$, $3x + 2y = 0$.
De la tercera, $y = -\frac{3x}{2}$; la primera da $z = x + y =
-\frac x2$; y compruébese en la segunda:
$2x - \frac{3x}{2} - \frac x2 = 0$: se cumple. Luego
$\ker u = \operatorname{Vect}\bigl((2, -3, -1)\bigr)$ (tomando
$x = 2$), de dimensión $1$.

Teorema del \omterm{def:b1:linmaps:rank}{rango}: $\operatorname{rk} u = 3 - 1 = 2$. \omterm{def:b1:linmaps:kerim}{Imagen}:
generada por las imágenes de la \omterm{def:b1:vspaces:free}{base} canónica, $u(e_1) = (1,2,3)$,
$u(e_2) = (1,1,2)$, $u(e_3) = (-1,1,0)$; las dos primeras son
\omterm{def:b1:vspaces:free}{libres} y el \omterm{def:b1:linmaps:rank}{rango} es $2$: base $\bigl((1,2,3), (1,1,2)\bigr)$.

No es \omterm{def:b1:logic:inj}{inyectiva} ($\ker \neq \{0\}$) ni \omterm{def:b1:logic:inj}{sobreyectiva} (\omterm{def:b1:linmaps:rank}{rango}
$2 < 3$): coherente con el~\cref{cor:b1:linmaps:samedim}.
\end{solution}

\begin{solution}{exo:b1:linmaps:3}
\emph{\omterm{def:b1:linmaps:kerim}{Núcleo}:} $P = P'$ fuerza $\deg P = \deg P'$ salvo que
$P = 0$; pero $\deg P' < \deg P$ para $P \neq 0$: luego
$\ker u = \{0\}$ y $u$, endomorfismo \omterm{def:b1:logic:inj}{inyectivo} del $\R_n[X]$ \omterm{def:b1:findim:def}{de
dimensión finita}, es un isomorfismo
(\cref{cor:b1:linmaps:samedim}).

\emph{Inversa:} sea $v(P) = P + P' + P'' + \dots + P^{(n)}$ (suma
finita en $\R_n[X]$). Entonces
\[
v\bigl(u(P)\bigr) = \sum_{k=0}^{n} (P - P')^{(k)}
= \sum_{k=0}^{n} P^{(k)} - \sum_{k=0}^{n} P^{(k+1)}
= P - P^{(n+1)} = P ,
\]
telescopando, ya que $P^{(n+1)} = 0$. Luego $v = u^{-1}$.
\end{solution}

\begin{solution}{exo:b1:linmaps:4}
$\operatorname{im}(v \circ u) = v(\operatorname{im} u) \subseteq
\operatorname{im} v$: \omterm{def:b1:linmaps:rank}{rango} $\leq \operatorname{rk} v$. Y $v$
restringida a $\operatorname{im} u$ tiene \omterm{def:b1:linmaps:kerim}{imagen}
$\operatorname{im}(vu)$; con el teorema del \omterm{def:b1:linmaps:rank}{rango} dentro de
$\operatorname{im} u$: $\operatorname{rk}(vu) \leq
\dim\operatorname{im} u = \operatorname{rk} u$.
\end{solution}

\begin{solution}{exo:b1:linmaps:5}
$u^2 = 0$ significa $u(u(x)) = 0$ para todo $x$: todo $u(x)$ está en
$\ker u$, es decir, $\operatorname{im} u \subseteq \ker u$. Y
entonces el teorema del \omterm{def:b1:linmaps:rank}{rango}:
\[
\dim E = \dim\ker u + \operatorname{rk} u \geq 2\operatorname{rk} u .
\]
Ejemplo con igualdad en $\R^2$: $u(x, y) = (y, 0)$: $u^2 = 0$ y
$\operatorname{rk} u = 1 = \frac{\dim E}{2}$.
\end{solution}

\begin{solution}{exo:b1:linmaps:6}
$(pq)^2 = pqpq = ppqq = pq$ (por la conmutación): una \omterm{def:b1:linmaps:projection}{proyección}
(\cref{thm:b1:linmaps:projchar}).

\omterm{def:b1:linmaps:kerim}{Imagen}: $\operatorname{im}(pq) \subseteq \operatorname{im} p$ y
$= \operatorname{im}(qp) \subseteq \operatorname{im} q$: está
contenida en la intersección. Recíprocamente, si
$x \in \operatorname{im} p \cap \operatorname{im} q$, entonces
$p(x) = x$ y $q(x) = x$ (los puntos fijos caracterizan la \omterm{def:b1:linmaps:kerim}{imagen} de
una \omterm{def:b1:linmaps:projection}{proyección}), luego $pq(x) = x$: $x \in \operatorname{im}(pq)$.

\omterm{def:b1:linmaps:kerim}{Núcleo}: $\ker p \subseteq \ker(qp) = \ker(pq)$ y, análogamente,
$\ker q \subseteq \ker(pq)$: la suma está contenida.
Recíprocamente, sea $pq(x) = 0$ y escríbase
\[
x = \underbrace{q(x)}_{\in\, \ker p} +
\underbrace{(x - q(x))}_{\in\, \ker q} :
\]
el primer término cumple $p(q(x)) = 0$, luego está en $\ker p$; y el
segundo está en $\ker q$, ya que $q(x - q(x)) = q(x) - q^2(x) = 0$.
Por tanto, $x \in \ker p + \ker q$.
\end{solution}

\begin{solution}{exo:b1:linmaps:7}
Nótense primero las inclusiones generales $\ker u \subseteq \ker u^2$
y $\operatorname{im} u^2 \subseteq \operatorname{im} u$ y, por el
teorema del \omterm{def:b1:linmaps:rank}{rango}, (2) $\iff$ (3) (\omterm{def:b1:linmaps:kerim}{núcleos} iguales $\iff$ \omterm{def:b1:linmaps:rank}{rangos}
iguales $\iff$ imágenes iguales, dadas las inclusiones).

(1 $\Rightarrow$ 2): sea $u^2(x) = 0$; entonces
$u(x) \in \ker u \cap \operatorname{im} u = \{0\}$, luego $x \in \ker u$.

(2 $\Rightarrow$ 1): por Grassmann y el teorema del \omterm{def:b1:linmaps:rank}{rango},
$\dim(\ker u + \operatorname{im} u) = \dim\ker u +
\operatorname{rk} u - \dim(\ker u \cap \operatorname{im} u) =
\dim E - \dim(\ker u \cap \operatorname{im} u)$: la suma es $E$ si y
solo si la intersección es $\{0\}$. Sea
$y \in \ker u \cap \operatorname{im} u$: $y = u(x)$ y $u(y) = 0$,
luego $u^2(x) = 0$ y, por (2), $u(x) = 0$: $y = 0$. Por tanto,
$E = \ker u \oplus \operatorname{im} u$.
\end{solution}

\begin{solution}{exo:b1:linmaps:8}
Si $\varphi = 0$: entonces $\ker\varphi = E \subseteq \ker\psi$
fuerza $\psi = 0 = 0\cdot\varphi$. Si $\varphi \neq 0$:
$\ker\varphi$ es un \omterm{def:b1:linmaps:forms}{hiperplano}; tómese $a \notin \ker\varphi$ y
póngase $\lambda = \frac{\psi(a)}{\varphi(a)}$. La forma
$\psi - \lambda\varphi$ se anula en $\ker\varphi$ (las dos lo hacen,
por la inclusión) y en $a$: se anula en
$\ker\varphi \oplus Ka = E$. Luego $\psi = \lambda\varphi$.
\end{solution}

\begin{solution}{exo:b1:linmaps:9}
Supóngase $\lambda_0 x + \lambda_1 u(x) + \dots + \lambda_{n-1}
u^{n-1}(x) = 0$. Aplíquese $u^{n-1}$: mueren todos los términos con
un factor $u^{\geq n}$ y queda $\lambda_0 u^{n-1}(x) = 0$, luego
$\lambda_0 = 0$. Aplíquese $u^{n-2}$ a la relación restante:
$\lambda_1 u^{n-1}(x) = 0$, luego $\lambda_1 = 0$; y así
sucesivamente. La familia es \omterm{def:b1:vspaces:free}{libre} y, al ser de tamaño
$n = \dim E$, es una \omterm{def:b1:vspaces:free}{base} (\cref{prop:b1:findim:twoofthree}). (En
esa \omterm{def:b1:vspaces:free}{base}, $u$ actúa como un desplazamiento --- el modelo de la
nilpotencia máxima.)
\end{solution}

\begin{solution}{exo:b1:linmaps:10}
\begin{enumerate}
  \item $f \circ f = -\mathrm{id}$ es \omterm{def:b1:logic:inj}{biyectiva}, luego $f$ lo es
        (\cref{prop:b1:logic:comp} adaptada: $f$ tiene la inversa
        bilátera $-f$). Si $f(x) = \lambda x$ con $x \neq 0$:
        aplicando $f$, $-x = \lambda^2 x$, luego $\lambda^2 = -1$:
        imposible en $\R$.
  \item Constrúyase la familia con voracidad. Tómese $x_1 \neq 0$:
        $(x_1, f(x_1))$ es \omterm{def:b1:vspaces:free}{libre} por (1). Si el
        $\operatorname{Vect}$ de la familia en curso
        $\bigl(x_1, f(x_1), \dots, x_k, f(x_k)\bigr)$, llamémoslo
        $V_k$ ---un \omterm{def:b1:vspaces:subspace}{subespacio} estable por $f$ (cada generador
        va a otro generador o a su opuesto:
        $f(f(x_i)) = -x_i$)---, no es todo $E$, elíjase
        $x_{k+1} \notin V_k$. \emph{Afirmación: la familia ampliada
        es \omterm{def:b1:vspaces:free}{libre}.} Supóngase
        $\alpha x_{k+1} + \beta f(x_{k+1}) + v = 0$ con $v \in V_k$ y
        $(\alpha, \beta) \neq (0,0)$. Aplíquese $f$:
        $\alpha f(x_{k+1}) - \beta x_{k+1} + f(v) = 0$ con
        $f(v) \in V_k$. Elimínese $f(x_{k+1})$ entre las dos
        relaciones (multiplíquese la primera por $\alpha$, la segunda
        por $-\beta$ y súmense):
        \[
        (\alpha^2 + \beta^2)\, x_{k+1} \in V_k ,
        \]
        y $\alpha^2 + \beta^2 \neq 0$ fuerza $x_{k+1} \in V_k$:
        contradicción. Así que la construcción continúa, añadiendo
        vectores \emph{de dos en dos}, hasta que $V_k = E$: la
        familia final es una \omterm{def:b1:vspaces:free}{base} de tamaño par, y $n$ es par.
\end{enumerate}
\end{solution}

\begin{solution}{exo:b1:linmaps:11}
\omterm{def:b1:reals:bounds}{Cota superior}: para todo $x$,
$(u + v)(x) = u(x) + v(x) \in \operatorname{im} u +
\operatorname{im} v$, luego
\[
\operatorname{rk}(u + v)
\leq \dim(\operatorname{im} u + \operatorname{im} v)
\leq \operatorname{rk} u + \operatorname{rk} v
\]
(Grassmann, \cref{thm:b1:findim:grassmann}). Cota inferior:
aplíquese la \omterm{def:b1:reals:bounds}{cota superior} a la pareja $(u + v, -v)$, cuya suma es $u$:
\[
\operatorname{rk} u \leq \operatorname{rk}(u + v) +
\operatorname{rk}(-v) = \operatorname{rk}(u + v) +
\operatorname{rk} v,
\]
luego $\operatorname{rk} u - \operatorname{rk} v \leq
\operatorname{rk}(u+v)$; e intercambiando $u$ y $v$ se obtiene el
valor absoluto.
\end{solution}

\begin{solution}{exo:b1:linmaps:12}
\emph{Fórmula exacta.} Sea $v'$ la restricción de $v$ al
\omterm{def:b1:vspaces:subspace}{subespacio} $\operatorname{im} w$. Su \omterm{def:b1:linmaps:kerim}{imagen} es
$v(w(F)) = \operatorname{im}(v \circ w)$ y su \omterm{def:b1:linmaps:kerim}{núcleo} es
$\ker v \cap \operatorname{im} w$. Teorema del \omterm{def:b1:linmaps:rank}{rango} para $v'$ sobre
el espacio $\operatorname{im} w$:
\[
\operatorname{rk} w = \dim\operatorname{im} w
= \operatorname{rk}(v \circ w) + \dim(\ker v \cap
\operatorname{im} w) .
\]

\emph{Frobenius.} Aplíquese dos veces la fórmula exacta, a $w$ y a
$w \circ u$:
\[
\operatorname{rk} w - \operatorname{rk}(vw)
= \dim\bigl(\ker v \cap \operatorname{im} w\bigr),
\qquad
\operatorname{rk}(wu) - \operatorname{rk}(vwu)
= \dim\bigl(\ker v \cap \operatorname{im}(wu)\bigr) .
\]
Como $\operatorname{im}(w \circ u) \subseteq \operatorname{im} w$, la
segunda intersección está contenida en la primera y su dimensión no
es mayor:
\[
\operatorname{rk}(wu) - \operatorname{rk}(vwu)
\;\leq\; \operatorname{rk} w - \operatorname{rk}(vw) ,
\]
que se reordena en la desigualdad de Frobenius. Con
$w = \mathrm{id}_F$ (de \omterm{def:b1:linmaps:rank}{rango} $\dim F$, y con
$\operatorname{im}\, \mathrm{id}_F = F$):
$\operatorname{rk} v + \operatorname{rk} u \leq \dim F +
\operatorname{rk}(vu)$, la desigualdad de Sylvester --- demostrada de
nuevo, en forma matricial, en el~\cref{exo:b1:matrices:10}.
\end{solution}

\begin{solution}{pb:b1:linmaps:1}
\textbf{1.} $(\mathrm{id} - p)^2 = \mathrm{id} - 2p + p^2 =
\mathrm{id} - p$: un proyector. Si $y = x - p(x)$, entonces
$p(y) = p(x) - p^2(x) = 0$ y, recíprocamente, $x \in \ker p$ da
$x = (\mathrm{id} - p)(x)$: $\operatorname{im}(\mathrm{id} - p) =
\ker p$. Y $(\mathrm{id} - p)(x) = 0 \iff p(x) = x \iff x \in
\operatorname{im} p$ (los puntos fijos,
\cref{thm:b1:linmaps:projchar}): $\ker(\mathrm{id} - p) =
\operatorname{im} p$.

\textbf{2.} $(\lambda\,\mathrm{id} + \mu p)^2 =
\lambda^2\,\mathrm{id} + (2\lambda\mu + \mu^2)\,p$. La pareja
$(\mathrm{id}, p)$ es \omterm{def:b1:vspaces:free}{libre} en $\mathcal{L}(E)$: $p = c\,
\mathrm{id}$ daría $c^2 = c$, luego $p = 0$ o $\mathrm{id}$,
excluidos. Identificando coeficientes, la \omterm{def:b1:logic:map}{aplicación} es un proyector
si y solo si $\lambda^2 = \lambda$ y $2\lambda\mu + \mu^2 = \mu$.
Para $\lambda = 0$: $\mu \in \{0, 1\}$. Para $\lambda = 1$:
$\mu^2 + \mu = 0$, $\mu \in \{0, -1\}$. Exactamente cuatro
proyectores en el plano: $0$, $p$, $\mathrm{id}$, $\mathrm{id} - p$.

\textbf{3.} $(\lambda\,\mathrm{id} + \mu p)(\lambda'\,\mathrm{id}
+ \mu' p) = \lambda\lambda'\,\mathrm{id} + (\lambda\mu' +
\mu\lambda' + \mu\mu')\,p$: el plano es estable por composición. Y
como $p^k = p$ para todo $k \geq 1$, para $Q = \sum_k a_k X^k$:
\[
Q(p) = a_0\,\mathrm{id} + \Bigl(\sum_{k \geq 1} a_k\Bigr) p
= Q(0)\,\mathrm{id} + \bigl(Q(1) - Q(0)\bigr)\,p .
\]

\textbf{4.} En $\operatorname{im} p$ (donde $p$ actúa como la
identidad), $\lambda\,\mathrm{id} + \mu p$ multiplica por
$\lambda + \mu$; y en $\ker p$, por $\lambda$. Como
$E = \operatorname{im} p \oplus \ker p$, la \omterm{def:b1:logic:map}{aplicación} es \omterm{def:b1:logic:inj}{biyectiva}
si y solo si $\lambda \neq 0$ y $\lambda + \mu \neq 0$. Resolviendo
$\lambda\alpha = 1$, $\lambda\beta + \mu\alpha + \mu\beta = 0$ en la
regla de composición de la pregunta 3:
\[
(\lambda\,\mathrm{id} + \mu p)^{-1}
= \frac1\lambda\,\mathrm{id} -
\frac{\mu}{\lambda(\lambda + \mu)}\,p ,
\]
cuya acción es por $1/\lambda$ en $\ker p$ y por
$1/(\lambda + \mu)$ en $\operatorname{im} p$, como debe ser.

\textbf{5.} Escríbase $F = \operatorname{im} p = \operatorname{im}
p'$. Para todo $x$, $p'(x) \in F$ y $p$ deja $F$ fijo punto a punto:
$p(p'(x)) = p'(x)$, es decir, $p\,p' = p'$; y simétricamente
$p'\,p = p$. Cuando dos proyecciones comparten su \omterm{def:b1:linmaps:kerim}{imagen}, decide la
que se aplica \emph{primero}: su salida ya está en $F$, donde la
\omterm{def:b1:linmaps:projection}{proyección} exterior actúa como la identidad y no cambia nada.

\textbf{6.} $(p + q)^2 = p^2 + pq + qp + q^2 = (p + q) + pq + qp$,
de modo que ser $p + q$ un proyector fuerza $pq + qp = 0$.
Compóngase por la izquierda con $p$: $pq + pqp = 0$; y por la
derecha con $p$: $pqp + qp = 0$. Restando, $pq = qp$; entonces
$pq + qp = 2pq = 0$ y la característica no es $2$: $pq = qp = 0$.

\textbf{7.} Con $pq = qp = 0$, el mismo desarrollo da
$(p + q)^2 = p + q$. \omterm{def:b1:linmaps:kerim}{Imagen}: siempre
$\operatorname{im}(p + q) \subseteq \operatorname{im} p +
\operatorname{im} q$. Recíprocamente, para $x \in \operatorname{im}
p$: $q(x) = q(p(x)) = 0$, luego $(p + q)(x) = p(x) = x$ y
$x \in \operatorname{im}(p+q)$; igual para $\operatorname{im} q$.
Que es directa: $x \in \operatorname{im} p \cap \operatorname{im} q$
da $x = p(x) = p(q(x)) = 0$. \omterm{def:b1:linmaps:kerim}{Núcleo}: si $p(x) + q(x) = 0$,
aplicar $p$ da $p(x) + p(q(x)) = p(x) = 0$, y aplicar $q$ da
$q(x) = 0$: $\ker(p + q) = \ker p \cap \ker q$ (la inclusión
recíproca es clara).

\textbf{8.} $p - q$ es un proyector si y solo si lo es
$\mathrm{id} - (p - q) = (\mathrm{id} - p) + q$ (pregunta 1, dos
veces). Por las preguntas 6--7 aplicadas a los proyectores
$\mathrm{id} - p$ y $q$, esto se cumple si y solo si
$(\mathrm{id} - p)q = q(\mathrm{id} - p) = 0$, es decir, si y solo
si $pq = q$ y $qp = q$.

\textbf{9.} $pq = q$ significa que $p$ deja fijo todo $q(x)$, es
decir, $\operatorname{im} q \subseteq \ker(p - \mathrm{id}) =
\operatorname{im} p$. Y $qp = q$ significa
$q\bigl((\mathrm{id} - p)(x)\bigr) = 0$ para todo $x$, es decir, que
$q$ se anula en $\operatorname{im}(\mathrm{id} - p) = \ker p$:
$\ker p \subseteq \ker q$. Los dos pasos son equivalencias: el orden
$q \leq p$ dice que $q$ proyecta sobre una \omterm{def:b1:linmaps:kerim}{imagen} menor,
paralelamente a un \omterm{def:b1:linmaps:kerim}{núcleo} mayor.

\textbf{10.} Desarrollando,
$(\mathrm{id} - p)(\mathrm{id} - q) = \mathrm{id} - p - q + pq =
\mathrm{id} - r$. Los proyectores $\mathrm{id} - p$ y
$\mathrm{id} - q$ conmutan, luego, por el~\cref{exo:b1:linmaps:6},
su producto $\mathrm{id} - r$ es el proyector sobre
$\operatorname{im}(\mathrm{id} - p) \cap
\operatorname{im}(\mathrm{id} - q) = \ker p \cap \ker q$
paralelamente a $\ker(\mathrm{id} - p) + \ker(\mathrm{id} - q) =
\operatorname{im} p + \operatorname{im} q$. Por la pregunta 1,
$r = \mathrm{id} - (\mathrm{id} - r)$ es entonces el proyector con
$\operatorname{im} r = \operatorname{im} p + \operatorname{im} q$ y
$\ker r = \ker p \cap \ker q$.

\textbf{11.} $p_i$ está bien definida (por la unicidad de la
descomposición) y es \omterm{def:b1:linmaps:def}{lineal} (la descomposición de $x + \lambda y$
es la suma de las descomposiciones, otra vez por la unicidad). Para
$x_i \in F_i$, la descomposición es el propio $x_i$, luego
$p_i(x_i) = x_i$: $p_i^2 = p_i$; y $p_j(x_i) = 0$ para $j \neq i$:
$p_i p_j = 0$ (pues $p_j(x) \in F_j$). Sumando las componentes,
$\sum_i p_i = \mathrm{id}$. Por último,
$\operatorname{im} p_i = F_i$ y $\ker p_i = \bigoplus_{j \neq i} F_j$.

\textbf{12.} $p_i = p_i \circ \mathrm{id} = p_i\sum_j p_j =
p_i^2 + \sum_{j \neq i} p_i p_j = p_i^2$: cada $p_i$ es un
proyector. Todo $x = \mathrm{id}(x) = \sum_i p_i(x)$ está en
$\sum_i \operatorname{im} p_i$: las imágenes suman $E$. Que es
directa: supóngase $y_1 + \dots + y_k = 0$ con
$y_i \in \operatorname{im} p_i$, de modo que $p_i(y_i) = y_i$.
Aplíquese $p_j$: $p_j (y_i) = p_j p_i (y_i) = 0$ para $i \neq j$,
luego $0 = p_j\bigl(\sum y_i\bigr) = y_j$, para todo $j$. Por tanto,
$E = \bigoplus_i \operatorname{im} p_i$.

\textbf{13.} $q = \mathrm{id} - p$, y la pregunta 1 da
$pq = p - p^2 = 0 = qp$ directamente: para dos proyectores, sumar la
identidad ya fuerza la ortogonalidad de la pareja.

\textbf{14.} $p + q = \mathrm{id} - r$ con $r$ proyector, y
$(\mathrm{id} - r)$ es un proyector (pregunta 1): luego $p + q$ es
un proyector, y la pregunta 6 da $pq = qp = 0$. Por \omterm{def:b1:linmaps:projection}{simetría}
($q + r = \mathrm{id} - p$ y $p + r = \mathrm{id} - q$), se anulan
todos los productos dos a dos, y la pregunta 12 concluye:
$E = \operatorname{im} p \oplus \operatorname{im} q \oplus
\operatorname{im} r$, automáticamente.

\textbf{15.} \emph{Lema.} Por inducción con Grassmann
(\cref{thm:b1:findim:grassmann}):
\[
\dim(F_1 + \dots + F_k) \leq \dim(F_1 + \dots + F_{k-1}) + \dim
F_k \leq \dots \leq \sum_i \dim F_i .
\]
Si el total es una igualdad, lo es cada paso:
$(F_1 + \dots + F_{j-1}) \cap F_j = \{0\}$ para todo $j$, y una
relación $y_1 + \dots + y_k = 0$ ($y_i \in F_i$) se derrumba desde la
derecha: $y_k \in (F_1 + \dots + F_{k-1}) \cap F_k = \{0\}$, después
$y_{k-1} = 0$, etc.: la suma es directa. Recíprocamente, una
\omterm{def:b1:vspaces:sum}{suma directa} tiene dimensiones aditivas (concaténense \omterm{def:b1:vspaces:free}{bases}).
\emph{Aplicación:} $x = \sum_i p_i(x)$ muestra
$E = \sum_i \operatorname{im} p_i$, luego
$n \leq \sum_i \operatorname{rk} p_i$; la hipótesis da la igualdad y,
por tanto, que es directa. Productos: fíjese $j$ y
$y \in \operatorname{im} p_j$. Entonces $y = \sum_i p_i(y)$ con
$p_i(y) \in \operatorname{im} p_i$, mientras que $y = y$ también es
una descomposición (solo la componente $j$); la unicidad fuerza
$p_i(y) = 0$ para $i \neq j$. Aplicado a $y = p_j(x)$: $p_i p_j = 0$.

\textbf{16.} Si $u^k(x) = 0$, entonces $u^{k+1}(x) = u(0) = 0$. Y
$\operatorname{im} u^{k+1} = u^k\bigl(u(E)\bigr) \subseteq
u^k(E) = \operatorname{im} u^k$.

\textbf{17.} Supóngase $\ker u^{r} = \ker u^{r+1}$ y sea
$x \in \ker u^{r+2}$: entonces $u(x) \in \ker u^{r+1} = \ker u^{r}$,
luego $u^{r+1}(x) = 0$: $x \in \ker u^{r+1}$. Con la pregunta 16,
$\ker u^{r+1} = \ker u^{r+2}$ y, por inducción, todos los \omterm{def:b1:linmaps:kerim}{núcleos}
posteriores coinciden con $\ker u^{r}$. Para las imágenes: el
teorema del \omterm{def:b1:linmaps:rank}{rango} da $\dim\operatorname{im} u^k = n - \dim\ker u^k$,
de modo que las dimensiones de las imágenes se congelan exactamente
cuando lo hacen las de los \omterm{def:b1:linmaps:kerim}{núcleos} y, con las inclusiones de la
pregunta 16, dimensiones iguales significan \omterm{def:b1:vspaces:subspace}{subespacios} iguales
(\cref{thm:b1:findim:subspaces}).

\textbf{18.} La sucesión $\bigl(\dim\ker u^k\bigr)_k$ es no
decreciente con valores en $\intint{0}{n}$; no puede crecer
estrictamente $n + 1$ veces, así que para algún $r \leq n$ se tiene
$\dim\ker u^{r} = \dim\ker u^{r+1}$ y, por tanto,
$\ker u^{r} = \ker u^{r+1}$ (inclusión más igualdad de dimensiones).
Tómese el $r$ mínimo; la pregunta 17 lo congela todo a partir de $r$,
imágenes incluidas.

\textbf{19.} Intersección: sea
$x \in \ker u^{r} \cap \operatorname{im} u^{r}$, digamos
$x = u^{r}(y)$ con $u^{r}(x) = 0$. Entonces $u^{2r}(y) = 0$, y
$\ker u^{2r} = \ker u^{r}$ (pregunta 17), luego $x = u^{r}(y) = 0$.
Dimensiones: el teorema del \omterm{def:b1:linmaps:rank}{rango} para $u^{r}$ da
$\dim\ker u^{r} + \dim\operatorname{im} u^{r} = n$; y con
intersección trivial, Grassmann hace de la suma un \omterm{def:b1:vspaces:subspace}{subespacio} de
dimensión $n$: $E = \ker u^{r} \oplus \operatorname{im} u^{r}$.

\textbf{20.} Estabilidad: $u^{r}(u(x)) = u(u^{r}(x)) = 0$ para
$x \in \ker u^{r}$; y
$u(u^{r}(y)) = u^{r}(u(y)) \in \operatorname{im} u^{r}$. En
$N = \ker u^{r}$: $(u|_N)^{r} = 0$ por definición de $N$: nilpotente.
En $I = \operatorname{im} u^{r}$:
$\ker(u|_I) = \ker u \cap I \subseteq \ker u^{r} \cap I = \{0\}$,
luego $u|_I$ es un endomorfismo \omterm{def:b1:logic:inj}{inyectivo} del $I$ \omterm{def:b1:findim:def}{de
dimensión finita} y, por tanto, \omterm{def:b1:logic:inj}{biyectivo}
(\cref{cor:b1:linmaps:samedim}).

\textbf{21.} Sea $x = a + b$ con $a \in N$, $b \in I$. Entonces
$u(x) = u(a) + u(b)$ con $u(a) \in N$ y $u(b) \in I$ (pregunta 20):
esa \emph{es} la descomposición de $u(x)$, luego
$\pi(u(x)) = u(a) = u(\pi(x))$: $\pi u = u\pi$.

\textbf{22.} $u^2(x,y,z) = u(y, 0, z) = (0, 0, z)$ y
$u^3(x,y,z) = u(0,0,z) = (0,0,z) = u^2(x,y,z)$. \omterm{def:b1:linmaps:kerim}{Núcleos}:
$\ker u = \{y = z = 0\} = \operatorname{Vect}(e_1)$,
$\ker u^2 = \{z = 0\} = \operatorname{Vect}(e_1, e_2)$,
$\ker u^3 = \ker u^2$: estabilización en $r = 2$. Imágenes:
$\operatorname{im} u = \operatorname{Vect}(e_1, e_3)$,
$\operatorname{im} u^2 = \operatorname{Vect}(e_3)$. Fitting:
$\R^3 = \operatorname{Vect}(e_1, e_2) \oplus
\operatorname{Vect}(e_3)$, y $\pi(x, y, z) = (x, y, 0)$.
Comprobación: $\pi u(x,y,z) = \pi(y, 0, z) = (y, 0, 0)$ y
$u\pi(x,y,z) = u(x, y, 0) = (y, 0, 0)$: iguales. En el primer factor,
$u(x, y, 0) = (y, 0, 0)$, cuyo cuadrado es $0$: nilpotente; y en el
segundo, $u(0,0,z) = (0,0,z)$: la identidad, \omterm{def:b1:logic:inj}{biyectiva}.

\textbf{23.} Si $u^m = 0$, entonces $\ker u^m = E$; y como los
\omterm{def:b1:linmaps:kerim}{núcleos} están congelados a partir de $r$,
$\ker u^{r} = \ker u^{\max(m, r)} = E$. Recíprocamente,
$\ker u^{r} = E$ significa $u^{r} = 0$. Y $\ker u^{r} = E \iff$ el
proyector de Fitting es sobre $E$ paralelamente a $\{0\}$, es decir,
$\pi = \mathrm{id}$. Por último, $r \leq n$ (pregunta 18) da: todo
endomorfismo nilpotente cumple $u^{n} = 0$ --- el índice de
nilpotencia nunca supera la dimensión.

\textbf{24.} Sea $m$ un índice de nilpotencia de $u|_A$:
$A \subseteq \ker u^{m} \subseteq \ker u^{\max(m,r)} = \ker u^{r}$.
Como $u|_B$ es \omterm{def:b1:logic:inj}{biyectiva},
$B = u(B) = u^{k}(B) \subseteq \operatorname{im} u^{k}$ para todo
$k$; en particular, $B \subseteq \operatorname{im} u^{r}$. Entonces
\[
n = \dim A + \dim B \leq \dim\ker u^{r} +
\dim\operatorname{im} u^{r} = n :
\]
las dos inclusiones son igualdades de dimensiones y, por tanto, de
\omterm{def:b1:vspaces:subspace}{subespacios}: $A = \ker u^{r}$, $B = \operatorname{im} u^{r}$.

\textbf{25.} (i) La parte III es un diccionario: las
descomposiciones $E = F_1 \oplus \dots \oplus F_k$ se corresponden
exactamente con las familias de proyectores con
$\sum p_i = \mathrm{id}$ y $p_i p_j = 0$, siendo los $F_i$ las
imágenes. (ii) Para $k = 3$, los complementarios
$\mathrm{id} - p_i$ son ellos mismos proyectores, lo que cerró el
argumento sin ninguna hipótesis extra; para $k$ general hace falta
$\sum_i \operatorname{rk} p_i \leq n$, una desigualdad que la traza
del segundo año da gratis ($\operatorname{tr} p = \operatorname{rk}
p$ para un proyector, y las trazas suman
$\operatorname{tr} \mathrm{id} = n$). (iii) El
\cref{exo:b1:linmaps:7} es el lema de Fitting en el caso ya
estabilizado $r \leq 1$; en general se deja que las cadenas de
\omterm{def:b1:linmaps:kerim}{núcleos} e imágenes se congelen, lo que lleva a lo sumo $n$ pasos.
(iv) En el volumen del segundo año, aplicado a
$u - \lambda\, \mathrm{id}$, el factor nilpotente se convierte en el
\omterm{def:b1:vspaces:subspace}{subespacio} propio generalizado en $\lambda$, y los proyectores de la
parte III, en los proyectores espectrales de la teoría de reducción.
El teorema de la parte IV es el \emph{lema de Fitting}.
\end{solution}
